%% =====================================================================
%% robo-chem -- IEEE Robotics and Automation Letters (RA-L) submission
%% =====================================================================
%% Build: latexmk -pdf main
%%        (or pdflatex main && bibtex main && pdflatex main && pdflatex main,
%%        or upload paper/ to Overleaf)
%%
%% RA-L rules this file follows (ieee-ras.org/publications/ra-l, checked
%% 2026-10-28 -- re-check before submitting):
%%  * Length: 6 pages in two-column format, EVERYTHING included (figures,
%%    tables, references, appendices). Up to 2 more pages at $175/page,
%%    8 pages absolute maximum. Per-section budgets are in the comment
%%    above each section: ~5.3 pages of body + ~0.7 page of references.
%%  * Initial submission: IEEE conference format, ieeeconf.cls from the
%%    PaperCept support pack (shipped next to this file), letter paper.
%%  * Final (accepted) version: IEEE journal format, IEEEtran.cls, with the
%%    RA-L running head and first-page footnotes. Flip \camerareadytrue below.
%%  * Double-anonymous review (all RAS journals since 2025-01-01): leave a
%%    byline but no names or affiliations; no acknowledgments or funding
%%    until accepted; cite your own work in the third person ("In [3]...",
%%    never "our earlier work"); no identifying links (GitHub, YouTube,
%%    lab pages) -- use https://anonymous.4open.science/ if needed; no lab
%%    names/logos/faces in figures or the video.
%%  * No cover page. Submit as PDF >= 1.4 with all fonts embedded; run it
%%    through the PaperCept PDF check before uploading.
%%  * Keywords: 2-5, the first two from the RA-L keyword list. They are
%%    also entered in PaperCept.
%%  * One optional video/multimedia zip (<= 50 MB, ReadMe.txt +
%%    Summary.txt). It may not carry extra text or results.
%% =====================================================================

% ---- Submission vs. final version -----------------------------------
% \camerareadyfalse: anonymous initial submission (ieeeconf, conference)
% \camerareadytrue:  accepted version (IEEEtran journal, names, footnotes)
\newif\ifcameraready
\camerareadyfalse

% \documentclass cannot sit inside an \if branch (it looks ahead past the
% \else), so the branch only picks the class and it is loaded after \fi.
%
% ieeeconf.cls is NOT part of TeX Live/Overleaf: it comes from the PaperCept
% pack and must sit next to main.tex (upload it to Overleaf too). If it is
% missing, the paper still builds with IEEEtran's conference mode, but page 1
% carries a red warning: that is not the format RA-L asks for.
\newif\ifieeeconfmissing
\ifcameraready
  \newcommand{\loadpaperclass}{%
    \documentclass[letterpaper, 10pt, journal, twoside]{IEEEtran}}
\else
  \IfFileExists{ieeeconf.cls}{%
    \newcommand{\loadpaperclass}{%
      \documentclass[letterpaper, 10pt, conference]{ieeeconf}%
      \IEEEoverridecommandlockouts  % enables \thanks and the keywords block
      \overrideIEEEmargins}         % required by PaperCept (printer margins)
  }{%
    \ieeeconfmissingtrue
    \newcommand{\loadpaperclass}{%
      \documentclass[letterpaper, 10pt, conference]{IEEEtran}%
      \IEEEoverridecommandlockouts}
  }
\fi
\loadpaperclass
\ifieeeconfmissing
  % IEEEtran calls ieeeconf's keywords environment IEEEkeywords.
  \newenvironment{keywords}{\begin{IEEEkeywords}}{\end{IEEEkeywords}}
  \typeout{*** ieeeconf.cls not found: built with IEEEtran conference mode.
    Put ieeeconf.cls next to main.tex before submitting to RA-L. ***}
\fi

% ---- Packages -------------------------------------------------------
\usepackage{cite}
\usepackage{amsmath,amssymb,amsfonts}
\usepackage{graphicx}
\usepackage{booktabs}
\usepackage{multirow}
\usepackage{makecell}      % two-line table headers
\usepackage{xcolor}
\usepackage{url}
\usepackage{balance}       % evens the columns on the last page
\usepackage[hidelinks]{hyperref}
% Keep author names out of the PDF metadata while under review.
\ifcameraready\else\hypersetup{pdfauthor={}, pdfcreator={}}\fi

\graphicspath{{figures/}}

% ---- Drafting helpers (delete before submission) --------------------
\newcommand{\todo}[1]{\textcolor{red}{[TODO: #1]}}
\newcommand{\note}[1]{\textcolor{blue}{[#1]}}
\newcommand{\tbd}{\textcolor{red}{--}}
% Framed stand-in so the paper compiles before figures exist.
% Usage: \placeholderfig{width}{height}{what goes here}
\newcommand{\placeholderfig}[3]{%
  \fbox{\parbox[c][#2][c]{#1}{\centering\footnotesize\textcolor{gray}{#3}}}}

% Shorthand used throughout
\newcommand{\sysname}{RoboChem}   % \todo{pick the system's name}
\newcommand{\eg}{e.g.,\ }
\newcommand{\ie}{i.e.,\ }

% Index terms: 2-5; the first two must come from the RA-L keyword list
% (https://www.ieee-ras.org/publications/ra-l/keywords).
\newcommand{\paperkeywords}{AI-Enabled Robotics, Task Planning, Perception
for Grasping and Manipulation, Robotics and Automation in Life Sciences}
% \todo{confirm against the RAS keyword list}

% ---- Running head (final version only) -------------------------------
% Wording required by RA-L; IEEEtran sets it in capitals.
\ifcameraready
\markboth{IEEE Robotics and Automation Letters. Preprint Version. Accepted Month, Year}
{Author \MakeLowercase{\textit{et al.}}: \sysname: Short Title}
\fi

\begin{document}
\bstctlcite{IEEEexample:BSTcontrol}


\ifcameraready
\title{\sysname: From a Printed Instruction Sheet to a Completed
Chemistry Experiment on a Robot Arm}
\else
\title{\LARGE \bf \sysname: From a Printed Instruction Sheet to a Completed
Chemistry Experiment on a Robot Arm}
\fi
% \todo{working title -- revisit once the results section is written}

\ifcameraready
  % First-page footnotes in the order RA-L lists them. The DOI line is
  % required verbatim. Fill the dates and editor from the acceptance email.
  \author{First Author$^{1}$, \todo{author list}, and Amir Barati Farimani$^{1}$%
  \thanks{Manuscript received: Month, Day, Year; Revised Month, Day, Year;
  Accepted Month, Day, Year.}%
  \thanks{This paper was recommended for publication by Editor
  Editor A.\ Name upon evaluation of the Associate Editor and Reviewers'
  comments.}%
  \thanks{This work was supported by \todo{funding}.}%
  \thanks{$^{1}$Department of Mechanical Engineering, Carnegie Mellon
  University, Pittsburgh, PA 15213, USA. \todo{check affiliations and emails}}%
  \thanks{Digital Object Identifier (DOI): see top of this page.}}
\else
  % RA-L: leave space for the byline, but no author information.
  \author{Anonymous Author(s)}
\fi

\maketitle
\ifcameraready\else
  \thispagestyle{empty}
  \pagestyle{empty}
\fi
\ifieeeconfmissing
  \noindent\todo{ieeeconf.cls not found, so this is IEEEtran conference
  format, not the RA-L submission format. Put ieeeconf.cls next to
  main.tex (on Overleaf: upload it to the project).}
\fi

% =====================================================================
% Abstract: 150-200 words, one paragraph, no citations.
% Positioning (agreed 2026-10-28): the robot does chemistry on the human's
% terms -- it adapts to a lab set up for people instead of the lab being
% instrumented for the robot. "Like a human" is always made concrete by
% the properties P1-P5 in the introduction; avoid "human reasoning",
% "thinks like a chemist" and "human-level".
% Shape: human lab -> automation inverts it -> what we built -> result.
% Index terms: set in \paperkeywords in the preamble.
% =====================================================================
\begin{abstract}
\todo{1-2 sentences: how people do a bench protocol -- a printed sheet,
containers labelled by hand and set down loosely, looking again before
each step.}
\todo{1 sentence: the contrast -- lab automation usually makes the lab
readable to the robot (fixed stations, fiducials/barcodes, scripts).}
\todo{2-3 sentences: what the system does, one clause per property --
reads the photographed protocol (P1), works on an ordinary bench (P2),
identifies containers by their written labels across cameras (P3),
re-perceives before every skill call (P4), sets how each skill runs,
including speed, from what is handled (P5).}
\todo{Result: one end-to-end number (runs, sub-tasks completed,
interventions, time) and one number each from the label and perturbation
experiments. Label number available now (Sec.~\ref{sec:exp_labels}): 160/160
locates found the right container on printed labels and the handwriting of
three people, and 0 of 96 returned the wrong container when it was moved
away from its label. Pick one; the first is the easier one to state in a
clause.}
\todo{Takeaway: one step closer to a robot doing bench chemistry the way
people do. Avoid "human reasoning", "thinks like a chemist",
"human-level".}
\end{abstract}

\ifcameraready
\begin{IEEEkeywords}\paperkeywords\end{IEEEkeywords}
\else
\begin{keywords}\paperkeywords\end{keywords}
\fi

% =====================================================================
% I. Introduction -- budget ~0.9 page (with the teaser figure)
% Paragraphs: (1) how a person does a bench protocol, (2) automation
% makes the lab robot-readable instead, (3) "on the human's terms" made
% concrete as P1-P5, (4) how the system realizes each,
% (5) contributions. P1-P5 are reused as the columns of Table I and as
% the structure of Sec. VII, so keep the labels stable.
% =====================================================================
\section{Introduction}
\label{sec:intro}
A student running a bench protocol works from a printed instruction sheet.
The reagents sit in containers that have been labelled by hand and set down
wherever there was room, and identical cups are told apart only by their
labels. The bench is shared, so a cup may not be where it was a minute ago;
the student looks again before each step and finds it. The student also
adjusts each step to what is being handled, carrying an empty cup quickly and
a full one slowly. None of this is arranged for a robot: the instructions,
the labels and the layout are all meant for people. 

Existing robots for automated experiments mostly adapt the experiment's
procedure, objects and layout to the system. The instructions have to be
written as a script or in a chemical programming
language~\cite{steiner2019organic,mehr2020universal,fakhruldeen2022archemist};
every container and instrument needs a fixed or known
position~\cite{burger2020mobile,dai2024autonomous}; and some systems
identify containers by fiducial markers such as AprilTags rather than by
the labels placed with them~\cite{yoshikawa2022chemistry}. This makes
automated experiments fast and accurate, but at the cost of flexibility: a
new experiment requires substantial manual reprogramming instead of simply
handing the system a printed instruction sheet, and if an object is moved
during the experiment, the run may not be able to continue.

This work aims to close this gap by letting a robotic system adapt to a lab
arranged for people. We make this goal concrete as five properties. Our system,   \sysname{}, (P1)~starts from the printed instruction sheet, with no manual
rewriting into a script; (P2)~works on an ordinary bench, where containers
need no fixed or preset positions and no visual markers; (P3)~tells
identical-looking containers apart by their hand-written or printed labels;
(P4)~looks at the bench again before every step, so it finds objects that
have been moved; and (P5)~adapts how each step is executed (how fast to move, how far to tip, how deep to dig) to what is handled and what the step asks for. Section~\ref{sec:related}
compares related work along these properties (Table~\ref{tab:related}), and
Section~\ref{sec:experiments} evaluates each of them.

\begin{figure}[t]
  \centering
  \placeholderfig{0.95\columnwidth}{4.2cm}{Teaser: the printed sheet
  (left) $\rightarrow$ the robot at a bench of loosely placed,
  hand-labelled cups (centre) $\rightarrow$ the finished experiment
  (right)}
  \caption{\todo{One sentence that states the framing on its own: the
  robot runs the worksheet the way a student would, from the printed
  sheet, on a bench of loosely placed, hand-labelled containers.}}
  \label{fig:teaser}
\end{figure}

We implement \sysname{} with all five properties on a Franka Panda arm
(Fig.~\ref{fig:teaser}). (P1)~A vision-language model (VLM) extracts the
steps from a photo of the sheet, and a planner turns them into parameterized
calls to a fixed set of skills. (P2)~SAM~3~\cite{carion2026sam3} segments
objects in four camera views, and positions come from multi-view geometry:
the depth point cloud for opaque objects and silhouettes for transparent
ones. The system only needs a list of the objects on the bench (with the
size of each transparent container), not their positions. (P3)~The VLM
reads each container's label in every view, and the cameras vote.
(P4)~Before each skill call, the target is located again if it has been
handled or its image has changed. (P5)~The VLM chooses each call's pour
angle from the wording of the step and its speed from what is held; a fixed
rule sets the angles when the same container is poured from, or into, more
than once, and the scoop depth comes from the measured powder surface.

The contributions of \sysname{} are as follows:
% Add each bullet's headline number once the experiments are run.
\begin{itemize}
  \item A system that carries out a chemistry experiment on a Franka Panda
  arm from a photo of the printed instruction sheet, on an ordinary bench
  with no fixed positions or markers for the containers.
  \todo{end-to-end result}
  \item A way to tell identical-looking containers apart by their
  hand-written or printed labels, which they stand on, read in each of four
  camera views and combined by a cross-camera vote. Because the position of
  each container comes from multi-view geometry, the label only needs to say
  what the container is and can be placed loosely. It found the right
  container in all 160 locates with printed labels and the handwriting of
  three people, and it never picked the wrong container in 96 locates with
  a container moved away from its label.
  \item Re-locating and re-identifying the containers before every step,
  so the robot finds a container a person has moved between steps, even
  among identical-looking ones. We compare this against locating objects
  once at the start. \todo{perturbation result}
\end{itemize}
% Bullet 3: if scan reuse (progress.md, 2026-10-22) is on in the evaluated
% code, say "before every skill call unless nothing has moved".


% =====================================================================
% II. Related Work -- budget ~0.6 page incl. Table I. Short paragraphs;
% end each with how this work differs. Table I carries the positioning:
% fill each row FROM THE PAPER ITSELF (only papers you have read and put
% in references.bib). Use \checkmark, a dash, or a short word ("partial",
% "sim only") -- do not guess.
% =====================================================================
\section{Related Work}
\label{sec:related}

\textbf{Robotic chemists and lab automation.}
Much of the reliability of automatic laboratories~\cite{tom2024self}
comes from adapting the experiment to the robot system. Procedures are written as scripts or in a chemical programming language~\cite{steiner2019organic,mehr2020universal,fakhruldeen2022archemist};
robots work between fixed stations, with fixed sample racks and taught
paths~\cite{dai2024autonomous,lunt2024modular}, and one localizes at each
station by touching a calibration cube to within
$\pm$0.12\,mm~\cite{burger2020mobile}; and containers or devices are
identified or localized with fiducial markers such as AprilTags or QR
codes~\cite{yoshikawa2022chemistry,angers2025roboculture,vescovi2023modular}.
More recently, large language models generate procedures from natural
language and pass them to liquid handlers or cloud
platforms~\cite{boiko2023autonomous,bran2024augmenting} or to robot arms at
the bench~\cite{yoshikawa2023large,darvish2025organa,wei2026agentchemist,zhang2025robochemist}. The closest to our system is
ORGANA~\cite{darvish2025organa}: it also uses a Franka arm and perceives
transparent glassware without markers, but the user assigns the role of each
identical-looking vessel at start-up, and it does not replan online. As
Table~\ref{tab:related} shows, none of these systems starts from a printed
instruction sheet, tells containers apart by reading their labels, or is
evaluated with a person moving objects mid-run.


\textbf{Language models for robot planning.}
The role of language models in robotics has grown from turning
instructions into skill sequences, programs or
trajectories~\cite{ahn2022saycan,singh2023progprompt,liang2023code,huang2023voxposer}
to setting execution parameters and closing the loop. For execution
parameters, Code as Policies~\cite{liang2023code} and
IROSA~\cite{knauer2026irosa} adjust speed in response to a user's
natural-language corrections, and ELLMER~\cite{monwilliams2025embodied} selects contact
forces from object properties. For closed-loop operation, Inner
Monologue~\cite{huang2022inner} re-detects the scene after every step, and
BioProVLA-Agent~\cite{du2026bioprovla} uses a VLM to check conditions before
and after each sub-task. Our system does not leave execution entirely to the
model: the model chooses the pour angle from the wording of the step and the
speed from what is held, a fixed rule assigns the angles when the same
container is poured from, or into, more than once, and the scoop depth comes
from the measured powder surface.

\textbf{Identifying objects in the lab.}
Fiducial markers such as AprilTags~\cite{olson2011apriltag,wang2016apriltag}
give both an exact identity and a pose, but every object needs its own tag,
and the offset from each tag to its object must be calibrated, which is
fragile and time-consuming to maintain~\cite{angers2025roboculture}.
Open-vocabulary models~\cite{liu2024grounding,carion2026sam3} need no
markers, but they find objects by category and cannot tell identical-looking
containers apart. Robots have identified objects by reading the text printed
on them with OCR~\cite{sharma2023semantic,liu2021robot,khan2025shake}. We
read the hand-written or printed label that each container stands on and
combine the reads across four views; the label only says what the container
is, while its position comes from multi-view geometry, so the same label can
be read by both people and the robot.


\begin{table}[t]
  \caption{Comparison along the properties of Sec.~\ref{sec:intro}:
  P1 instruction input, P2 no fixed positions or markers, P3 container
  identification, P4 re-perception before each step, P5 adapted execution.
  \checkmark: yes; $\circ$: partial; --: no or not reported.}
  \label{tab:related}
  \centering
  \footnotesize
  \begin{tabular}{@{}lcccccc@{}}
    \toprule
    System & P1 & P2 & P3 & P4 & P5 & Platform \\
    \midrule
    Chemputer~\cite{steiner2019organic,mehr2020universal} & script & -- & -- & -- & -- & Fluidic \\
    Burger et al.~\cite{burger2020mobile}       & script & -- & --      & --      & --      & Mobile arm \\
    Coscientist~\cite{boiko2023autonomous}      & text   & -- & --      & --      & --      & Liquid handler \\
    CLAIRify~\cite{yoshikawa2023large}          & text   & $\circ$ & markers & $\circ$ & --      & Arm \\
    ORGANA~\cite{darvish2025organa}             & text   & $\circ$ & user    & --      & --      & Arm \\
    AgentChemist~\cite{wei2026agentchemist}     & text   & $\circ$ & --      & $\circ$ & $\circ$ & Dual-arm \\
    RoboChemist~\cite{zhang2025robochemist}     & text   & $\circ$ & --      & $\circ$ & $\circ$ & Dual-arm \\
    \midrule
    \sysname{} (ours)                           & photo  & \checkmark & labels & \checkmark & \checkmark & Arm \\
    \bottomrule
  \end{tabular}
  % Fill our row only with what the experiments show. P5 (planner-set
  % speed) had not run on the robot as of 2026-10-27, and P4 depends on
  % whether scan reuse is on in the evaluated code version.
\end{table}

% =====================================================================
% III. System Overview -- budget ~0.4 page incl. the architecture figure.
% Hardware, software split, the data flow in Fig. 2. Keep it short: the
% next three sections carry the detail.
% =====================================================================
\section{System Overview}
\label{sec:overview}

\textbf{Hardware.} 

The system runs on a Franka Emika Panda arm with a Franka Hand
parallel-jaw gripper. Four Intel RealSense D415 RGB-D cameras are mounted
on a cage around the bench and view it from different sides; their
extrinsics are calibrated to the robot base with a cube held by the gripper,
which is used only for calibration. The bench holds ordinary lab equipment:
shallow white dishes of powdered reagents, clear plastic cups, a plastic
beaker, a plastic transfer pipette standing in a lidded paper cup, and a
3D-printed scoop, stirrer, and stirrer holder. Each container is stacked on a
slip of paper with its contents written or printed on it.

\textbf{Software.} 
The system runs as two processes: the segmentation model
SAM~3~\cite{carion2026sam3} runs as a separate GPU service, which the
program controlling the arm calls over the local network. The agents follow
the modular structure of PLATO~\cite{car2026plato} and consist of four
parts. An instruction parser extracts the goal, steps, reagents and
equipment from the photo of the sheet. A scene agent matches the objects
seen by the four cameras to the names in the list of bench objects. A
planner decomposes the goal into ordered sub-tasks and replans when a step
fails. Before each sub-task, a skill-call agent produces one skill call from
the measured state of the gripper; its output is constrained by a JSON
schema to a fixed set of skills. The scene agent uses GPT-4.1 mini and the
other three use GPT-4.1. \todo{probably needs to be changed}
Each manipulation skill performs its own perception, grasp
analysis and motion, and checks during execution that it has reached its
target.

\textbf{Data flow.} \todo{Walk through Fig.~\ref{fig:architecture} in 4-5
sentences.}

\begin{figure*}[t]
  \centering
  \placeholderfig{0.95\textwidth}{5.0cm}{Architecture: instruction image
  $\rightarrow$ scene agent $\rightarrow$ planner $\rightarrow$ skill-call
  agent $\rightarrow$ skills (perception + motion) $\rightarrow$ checks
  $\rightarrow$ report; sim twin shown alongside}
  \caption{\todo{System architecture.}}
  \label{fig:architecture}
\end{figure*}

% =====================================================================
% IV. Perception -- budget ~0.7 page.
% Source material: progress.md "Calibration / perception notes",
% the 2026-10-10 recalibration entries, label voting (2026-10-20).
% =====================================================================
\section{Multi-View Perception}
\label{sec:perception}

\subsection{Open-Vocabulary Segmentation and 3D Fusion}

At every object identification request, SAM is run on each of the four cameras and ranks objects it sees to closeness of the object being searched for. The top 2 camera's scans which are closest to the object prompted are kept as the truth for pointcloud identification to go off of, because other cameras might not have a view of the label or object. The points identified are then bak-projected onto the pointcloud scan to get the 3d location of where the gripper arm should go to in order to complete the requested action on the requested object. 

\subsection{Reading Human-Written Labels (P3)}

Human labels are placed adjacent to objects that are ambiguous to humans
as well

\todo{Identical clear cups told apart by reading their labels in each view
with the VLM, then a vote across cameras; ties resolved by the last known
position. Make the trade-off explicit: pose comes from segmentation and
the fused cloud, so the label only has to carry identity, and text does
that with no setup, no placement constraint, and is readable by people.
Do not claim it beats fiducials at what fiducials do (exact ID and pose).}

\todo{The experiments in Sec.~\ref{sec:exp_labels} report locates found
"by label" and "by elimination", so this subsection has to define both
first. (1)~Each camera's crop of each candidate is read by the VLM as a
closed-set choice among the label strings on the bench, not free OCR.
(2)~Containers of the same kind form a vote group (the two white dishes;
the clear cup and the beaker). (3)~Elimination: if no camera reads the
wanted label but every other label in the group is placed firmly (two or
more cameras) on another container, the one container left is the
answer (\texttt{robochem/vision/\_\_init\_\_.py}, \texttt{eliminate}). (4)~If
neither works the locate returns "not found" rather than a guess; say so
here, because the distance test shows this is why misses are never wrong
picks.}

\todo{Say where the label must be for a read to work, from the distance
test: the label is read when it is inside the container's crop, i.e.\ the
container stands on or right against its slip of paper. The user's
sentence above says "adjacent"; keep that, and make it match what was
measured.}

\begin{figure}[t]
  \centering
  \placeholderfig{0.95\columnwidth}{3.5cm}{Four camera views with masks
  and the fused point cloud}
  \caption{\todo{Perception pipeline output on the real bench.}}
  \label{fig:perception}
\end{figure}

% =====================================================================
% V. Manipulation Skills -- budget ~0.9 page incl. Table I.
% One short paragraph per skill family; put parameters in Table I rather
% than in prose. Say how each skill detects its own failure.
% =====================================================================
\section{Manipulation Skills}
\label{sec:skills}

\todo{Common structure: scan (P4) $\rightarrow$ localize $\rightarrow$
plan motion $\rightarrow$ execute $\rightarrow$ check. Execution adapted
to what is handled (P5): every call carries a speed (0--100) chosen by the
planner, where 0 is the slowest, careful pace used for liquids and 100 is
2.5$\times$ faster (robochem/skills/pacing.py), alongside pour angle and
scoop depth.}

\textbf{Grasping.} \todo{Upright top grasp, width = measured diameter
minus squeeze; yaw alignment for tools lying at an angle.}

\textbf{Scoop and dump.} \todo{Stroke chosen from the powder-height map;
force-guarded descent; dump with shake and wrist twist; sideways headings
when reach limits the tip angle.}

\textbf{Pour.} \todo{Lip of held cup placed over the target; tip-angle
tracking and stall detection.}

\textbf{Stir and tool return.} \todo{Circular stir path inside the cup;
stirrer re-seated in its holder; scoop returned to its pick site.}

\begin{table}[t]
  \caption{Skill library. \todo{fill}}
  \label{tab:skills}
  \centering
  \footnotesize
  \begin{tabular}{@{}lll@{}}
    \toprule
    Skill & Key parameters & Failure check \\
    \midrule
    \texttt{pick\_up} & \tbd & \tbd \\
    \texttt{scoop}    & \tbd & \tbd \\
    \texttt{dump}     & \tbd & \tbd \\
    \texttt{pour}     & \tbd & \tbd \\
    \texttt{stir}     & \tbd & \tbd \\
    \texttt{place}    & \tbd & \tbd \\
    \bottomrule
  \end{tabular}
\end{table}

% =====================================================================
% VI. Instruction Following and Planning -- budget ~0.5 page.
% Source: robochem/agents/, robochem/orchestrator/agent_orchestrator.py,
% robochem/agents/reporter.py. Name the model(s) and version used for the
% reported runs.
% =====================================================================
\section{From Instruction Sheet to Skill Calls}
\label{sec:planning}

\textbf{Instruction parsing.} \todo{Photographed sheet $\rightarrow$
goal and ordered steps.}

\textbf{Scene grounding.} \todo{Scene agent maps the bench photos to the
inventory; reagent-named containers.}

\textbf{Planning and skill calls.} \todo{Planner emits sub-tasks; the
skill-call agent fills parameters under a strict JSON schema over a closed
skill vocabulary; the executor fills tool geometry from CAD and carries
measured offsets between steps.}

\textbf{Execution, recovery and reporting.} \todo{Preflight reach check,
step-through mode, resume from a record, photos after each pour, and the
experiment report written after the run.}

% =====================================================================
% VII. Experiments -- budget ~1.2 pages. This is the section reviewers
% read most closely. State the protocol, the number of trials, and what
% counts as success BEFORE giving numbers.
%
% Data that already exists (see progress.md; check before quoting):
%   - 2026-10-08 full Magic Beaker run: 31/31 sub-tasks, 2177 s, no replans
%     (experiments/agent_run_20261008_221001.json)
%   - 2026-10-15 / 10-19 / 10-20 runs with their failures and resumes
%   - 2026-10-22 four-item jellyfish chain timings
% Most of these were single runs during development. A results table for
% the paper needs repeated trials with a fixed code version.
%
% Label experiments, 2026-10-29 (cameras only, arm never moved). These ARE
% repeated trials on one code version and go into Sec. VII-C as they are:
%   - Label reading benchmark, printed vs handwritten
%     (scripts/label_reading_benchmark.py,
%     experiments/label_benchmark_20261029/ on the robot PC only)
%   - Label distance test, labelled vs unlabelled
%     (scripts/label_distance_test.py,
%     experiments/label_distance_20261029/ on the robot PC only)
% The raw data (scans.jsonl, report.md) is gitignored and lives on the
% robot PC only. Every number below was checked against report.md and
% scans.jsonl there on 2026-11-04; the per-camera counts were computed
% from label_matches in scans.jsonl (each camera's read on the container
% the vote chose).
% =====================================================================
\section{Experiments}
\label{sec:experiments}

% Structure follows P1-P5: one subsection per claim, each with a baseline.
% The human-lab framing commits us to a messy setup: loosely placed and
% handwritten labels, and a person perturbing the bench. Trials on a
% carefully arranged bench undercut the claim.

\subsection{Setup and Protocols}
\todo{Tasks (\eg Magic Beaker, jellyfish, single-skill trials), number of
trials, how the bench is set up (labels placed by hand, layouts varied),
models used, success definitions, and the code version.}

\subsection{End-to-End from the Printed Protocol (P1, P2)}
\todo{Full runs from the photographed sheet: sub-tasks completed, replans,
human interventions, wall-clock time, final outcome (colour change).
Table~\ref{tab:e2e}.}

\subsection{Identifying Containers by Their Labels (P3)}
\label{sec:exp_labels}

Both label experiments use the same bench. Four containers stand in a
square in front of the arm, each on its paper label: two identical white
dishes (red cabbage powder and baking soda), a clear plastic cup (10\,ml
water) and a plastic beaker (vegetable oil). The reader chooses among
exactly these four label strings. Only the cameras are used; the arm does
not move. Each locate is the same call a skill makes, and every memory is
cleared before it (kept scans, last known position and tracked objects),
so no locate is helped by the one before. A locate is settled \emph{by
label} when at least one camera reads the wanted label on the container
the vote chooses, and \emph{by elimination} when no camera reads it but
every other label of that kind of container is placed on another one.
Otherwise it returns \emph{not found} instead of a guess. The ground truth
is each container's position located with label reading switched off and
named by its corner of the square; three repeats agreed within 0.4\,mm. A
locate with all four cameras took 14.1\,s on average (13.2 to 18.1\,s).

\textbf{Label legibility.} The labels were swapped between four sets:
printed, and handwritten by three different people. In each set every
container was located 10 times (Table~\ref{tab:label_reading}). All 160
locates found the right container, with a mean offset of 0.2\,mm from the
ground truth and at most 1.1\,mm. The vote across cameras is needed for
this. In 21 of the 160 locates one camera read another label on the right
container: camera 3 in 11, always on the red cabbage dish and on both
printed and handwritten paper, probably because it sees that paper
sideways at a low angle; camera 4 in 9, all on the beaker and 8 of them in
one handwriting set; and camera 2 once. No locate had more than one camera
misreading, so the other three outvoted it every time. Used alone, camera
3 would have found at most 149 of the 160 containers
(Table~\ref{tab:ablations}).

\begin{table}[t]
  \caption{Label reading on four label sets, each container located 10
  times per set. Offsets are from the ground truth located without
  labels. The last column counts locates in which all four cameras read
  the right label on the container.}
  \label{tab:label_reading}
  \centering
  \footnotesize
  \begin{tabular}{@{}lcccc@{}}
    \toprule
    Label set & Locates & Correct & \makecell{Offset (mm)\\mean / max} & \makecell{All 4 cameras\\read it} \\
    \midrule
    Printed       & 40  & 40/40   & 0.2 / 0.4 & 35/40 \\
    Handwriting 1 & 40  & 40/40   & 0.2 / 1.1 & 33/40 \\
    Handwriting 2 & 40  & 40/40   & 0.2 / 0.5 & 32/40 \\
    Handwriting 3 & 40  & 40/40   & 0.2 / 0.8 & 39/40 \\
    \midrule
    Total         & 160 & 160/160 & 0.2 / 1.1 & 139/160 \\
    \bottomrule
  \end{tabular}
  % From experiments/label_benchmark_20261029/ on the robot PC (report.md
  % and scans.jsonl, checked 2026-11-04). Sets in the order scanned on
  % 2026-10-29: printed, then three writers (names in scans.jsonl; keep
  % them out of the paper while under review). If you mean to call any set
  % "bad handwriting" in the text, put an example crop of each set in a
  % figure so the reader can judge.
\end{table}

\textbf{Distance from the label.} One container at a time (the baking
soda dish, the water cup or the beaker) was moved away from its label to
5 or 6 spots between 2 and 44\,cm, and located 3 times at each spot, 96
locates in all. The distance is how far its centre moved, measured by
hand. Each container was tested twice: with the other three labels left on
the bench, and with those labels removed while their containers stayed
(Table~\ref{tab:label_distance}). None of the 96 locates returned the
wrong container; every miss was \emph{not found}. With the other labels
present no locate failed, because whenever the moved container's own
label was not read, elimination found it (10 locates). With the other
labels removed, 16 of 48 locates were not found. The distance of the
centre does not predict whether the label is read. The water cup was read
by all four cameras 18\,cm from its spot, standing against the edge of its
paper, and was not found at 16\,cm with bare table between it and the
paper. The baking soda dish was not read at 6\,cm, where it stood on top
of its own text. What decides the read is whether the label is still in
view beside the container, that is, whether the container stands on or
against its paper. We saw this in a few frames and did not measure the gap
from the container to its paper or the direction of the move.

\begin{table}[t]
  \caption{Locates with one container moved away from its label, 3 per
  spot. ``Others' labels'' says whether the other three labels stayed on
  the bench; ``wrong'' is a locate that returned another container.}
  \label{tab:label_distance}
  \centering
  \footnotesize
  \setlength{\tabcolsep}{4pt}
  \begin{tabular}{@{}llccccc@{}}
    \toprule
    Moved container & \makecell{Others'\\labels} & Locates & \makecell{By\\label} & \makecell{By\\elim.} & Wrong & \makecell{Not\\found} \\
    \midrule
    \multirow{2}{*}{Baking soda dish} & on  & 18 & 14 & 4 & 0 & 0 \\
                                      & off & 18 &  9 & 1 & 0 & 8 \\
    \midrule
    \multirow{2}{*}{Water cup}        & on  & 15 & 12 & 3 & 0 & 0 \\
                                      & off & 15 & 12 & 0 & 0 & 3 \\
    \midrule
    \multirow{2}{*}{Oil beaker}       & on  & 15 & 12 & 3 & 0 & 0 \\
                                      & off & 15 &  9 & 1 & 0 & 5 \\
    \midrule
    Total                             &     & 96 & 68 & 12 & 0 & 16 \\
    \bottomrule
  \end{tabular}
  % From experiments/label_distance_20261029/report.md on the robot PC
  % (checked 2026-11-04; the 3 locates with no label reads at all, a
  % network failure, were on-label checks and are not in this table).
  % Spots ran from
  % about 6 cm to 44 cm; the farthest spot that was always read by label
  % was 11-21 cm depending on the container. If there is room, a column
  % "farthest always read" (13.1, 18, 12.3, 18, 20.8, 11.2 cm) helps, but
  % keep the caveat that direction from the paper was not recorded.
\end{table}

Two results need care. First, with the other labels removed, elimination
still settled 2 of the 48 locates, because the reader named a container
from what it holds: one camera read the red cabbage dish as ``red cabbage
powder'' from its red powder, and another read the water cup as ``10\,ml
water''. The rows without the other labels are therefore not a clean
condition with no other labels. Second, two far reads counted as by label
rest on a single camera and were not checked against the frames: the
baking soda dish at 44\,cm (2 of 3 locates) and at 18\,cm without the
other labels. They may come from the contents of the dish and not from its
label. Elimination was not switched off in the code for the phase without
the other labels. Removing those labels takes away the moved container's
only elimination partner without changing the system, so the test runs the
same code as a real run.

\subsection{Robustness to People Moving Things (P4)}
\todo{A person moves or swaps a container between steps at random.
Ablation: re-perceive before every skill vs.\ locate once at the start
(Table~\ref{tab:ablations}). Also report the cost: full scans took
11.6--15.3\,s per skill in the 2026-10-22 chain (progress.md).}

\subsection{Execution Adapted to What Is Handled (P5)}
\todo{Whole-protocol time and spills or failures with planner-set speed
vs.\ fixed slow (speed 0) vs.\ fixed fast (Table~\ref{tab:ablations}).
Not yet run on the robot as of 2026-10-27.}

\begin{table}[t]
  \caption{Ablations, one per property. \todo{fill}}
  \label{tab:ablations}
  \centering
  \footnotesize
  \begin{tabular}{@{}llcc@{}}
    \toprule
    Property & Condition & Success & Time (s) \\
    \midrule
    \multirow{2}{*}{P3 labels}  & Single camera (cam.~3) & $\leq$149/160 & -- \\
                                & Cross-camera vote    & 160/160 & 14.1 \\
    % P3: the vote row is the label reading benchmark (all four sets,
    % 14.1 s mean per locate). Single camera: each camera's read on the
    % container the vote chose (label_matches in scans.jsonl): cam 2
    % 159/160, cam 3 149/160, cam 4 151/160, cam 5 160/160. It is an upper
    % bound: a camera alone also fails when it reads the wanted label on a
    % second instance, and the jsonl cannot tell a duplicate mask of the
    % same container from a neighbour (no per-instance positions logged).
    % Not timed alone, hence "--".
    \midrule
    \multirow{2}{*}{P4 rescan}  & Locate once at start & \tbd & \tbd \\
                                & Before every skill   & \tbd & \tbd \\
    \midrule
    \multirow{3}{*}{P5 speed}   & Fixed slow           & \tbd & \tbd \\
                                & Fixed fast           & \tbd & \tbd \\
                                & Planner-set          & \tbd & \tbd \\
    \bottomrule
  \end{tabular}
\end{table}

\subsection{Skill-Level Results}
\todo{Keep short. Repeated trials per skill: success rate, time, and a
measured quantity (\eg powder mass per scoop, tip angle reached, poured
volume).}

\begin{table}[t]
  \caption{Per-skill results on the real robot. \todo{fill}}
  \label{tab:skill_results}
  \centering
  \footnotesize
  \begin{tabular}{@{}lccc@{}}
    \toprule
    Skill & Success & Time (s) & Measured quantity \\
    \midrule
    Pick (beaker/cup) & \tbd/\tbd & \tbd & \tbd \\
    Pick (scoop)      & \tbd/\tbd & \tbd & \tbd \\
    Scoop             & \tbd/\tbd & \tbd & \tbd \\
    Dump              & \tbd/\tbd & \tbd & \tbd \\
    Pour              & \tbd/\tbd & \tbd & \tbd \\
    Stir              & \tbd/\tbd & \tbd & \tbd \\
    Tool return       & \tbd/\tbd & \tbd & \tbd \\
    \bottomrule
  \end{tabular}
\end{table}

\begin{table}[t]
  \caption{End-to-end runs from the photographed protocol. \todo{fill}}
  \label{tab:e2e}
  \centering
  \footnotesize
  \begin{tabular}{@{}lcccc@{}}
    \toprule
    Protocol & Runs & Sub-tasks done & Interventions & Time (min) \\
    \midrule
    Magic Beaker & \tbd & \tbd & \tbd & \tbd \\
    Jellyfish    & \tbd & \tbd & \tbd & \tbd \\
    \bottomrule
  \end{tabular}
\end{table}

\subsection{Failure Analysis}
\todo{Where runs stopped and why: reach limits on far containers, label
ties, network/API errors, scoop jams. A small breakdown figure or table.}

\begin{figure}[t]
  \centering
  \placeholderfig{0.95\columnwidth}{3.5cm}{Key-frame strip of one full run}
  \caption{\todo{Key frames from an end-to-end run.}}
  \label{fig:keyframes}
\end{figure}

% =====================================================================
% VIII. Limitations -- budget ~0.25 page. Be specific; reviewers trust a
% paper more when it names its own failure modes.
% =====================================================================
\section{Limitations}
\label{sec:limitations}

\todo{Where the robot still falls short of a person. Measuring amounts:
no closed-loop measurement of dispensed mass. Reach limits at the far edge of the
bench. Occlusion of the scoop bowl underside (measured depth is a lower
bound). Reliance on a hosted VLM. One bench and a small set of
protocols.}

\todo{Labels, from Sec.~\ref{sec:exp_labels}: a label is only read while
the container stands on or against it, so a container a person puts
down away from its paper is found only by elimination, or not at all if
the other labels are gone too. When it fails, it reports "not found"
rather than picking the wrong container. Each locate takes about 14\,s
and needs the hosted VLM; with the network down, a locate took 134\,s and
read nothing. The labels tested were 4 short strings on one bench.}

% =====================================================================
% IX. Conclusion -- budget ~0.1 page. 3-4 sentences, no new claims.
% =====================================================================
\section{Conclusion}
\label{sec:conclusion}

\todo{A robot that does bench chemistry on the human's terms (P1--P5),
the headline result, and the next step.}

% Acknowledgments and funding are added only after acceptance.
\ifcameraready
\section*{Acknowledgment}
\todo{Lab members, funding, hardware.}
\fi

\balance
\bibliographystyle{IEEEtran}
\bibliography{references}

\end{document}
